\phantomsection
\addcontentsline{toc}{chapter}{General Introduction}
\pagenumbering{arabic}
\setcounter{page}{1}
{\huge\bfseries General Introduction}
\markboth{General Introduction}{}   % <-- add this

\vspace{1cm}

\initial{T}he Fourth Industrial Revolution, or Industry~4.0, has fundamentally
transformed industrial operations by integrating artificial intelligence~(AI),
the Internet of Things~(IoT), cloud computing, and advanced automation into
manufacturing and operational processes. As industrial systems become
increasingly digitized and interconnected, engineers and technicians require
secure, remote access to specialized equipment - programmable logic
controllers~(PLCs), human-machine interfaces~(HMIs), SCADA systems, and
industrial robots - for training, development, and maintenance purposes.
Traditional approaches to industrial remote access present significant
limitations: VPN-based solutions are complex to configure and, in many
Operational Technology~(OT) deployment contexts, incompatible with OT security
requirements; conventional remote desktop tools lack native integration with
industrial communication protocols; and no standardized platform has emerged to
consolidate virtual machine~(VM) provisioning, browser-based remote desktop
access, and IoT device control within a single, coherent system. As the
comparative analysis presented in Chapter~\ref{chap:preliminary_study} will demonstrate, existing solutions
address these challenges only partially. This raises a central engineering
question: \textit{how can industrial engineers and trainers securely access
remote equipment from any browser, without client-side software installation,
while preserving OT network integrity and enforcing strict access control?}

It is within this industrial and technological context that the present
internship was conducted at \textbf{Novation City} in Sousse, Tunisia, from
FEB~2026 to MAY~2026, as part of the final-year project requirements
for the Software Engineering Bachelor at EPI Digital School. Novation City is a competitiveness
cluster established in 2009 as a public-private partnership dedicated to
fostering innovation and industrial transformation within the mechatronics and
Industry~4.0 sectors~{\footnotesize\cite{novationcity2026}}. Novation City operates through
several specialized business units, among which \textbf{Novation Akademy}
(industrial training and professional skills development) and \textbf{Innovation}
(applied research and R\&D initiatives) are the primary drivers of the
operational requirement for a secure remote laboratory infrastructure. This
convergence of training needs and R\&D capability motivated the development of
a dedicated platform to provide trainees and engineers with secure remote
access to industrial equipment without the constraints of physical presence.

To address these challenges, Novation City initiated the
\textbf{IoT-REAP} (\textbf{Internet of Things - Remote Equipment Access
Platform}) project. The platform provides industrial engineers, technicians, and
trainers with browser-based access to virtual machines through Proxmox~VE
integration~{\footnotesize\cite{proxmox-ve-api}}, and clientless remote desktop sessions via
Apache Guacamole, supporting the RDP, VNC, and SSH
protocols~{\footnotesize\cite{guacamole-manual}}. Camera controls are dispatched and monitored
via the lightweight MQTT messaging protocol~{\footnotesize\cite{mqtt-oasis}}, enabling
real-time communication with connected industrial devices. IoT-REAP adheres to
strong industrial cybersecurity principles, implementing zero-trust access
controls and enforcing OT network segmentation drawing upon the principles of
the Purdue reference model. The platform requires no client-side software
installation, thereby reducing the client-side attack surface and simplifying
access provisioning for geographically distributed users.

The primary objective of this internship was therefore to design and implement
the IoT-REAP platform, with specific focus on the development of the unified
web portal, the integration of Proxmox~VE and Apache Guacamole for
virtualization and remote desktop orchestration, and the implementation of the
secure MQTT-based IoT control interface. The deployment of underlying network
infrastructure, the administration of the Proxmox hypervisor layer, and the
development of proprietary hardware drivers lie outside the scope of this report
and were addressed by Novation City's internal IT operations team as part of the
broader infrastructure initiative.

The report comprises three chapters documenting the complete development
lifecycle of IoT-REAP. \textbf{Chapter~\ref{chap:preliminary_study}} establishes the professional context
at Novation City, characterizes the challenges of secure industrial remote
access, critically reviews existing industrial remote access solutions and their
limitations, defines the platform requirements and scope, and presents the Agile
Scrum methodology adopted, including the five-sprint development timeline.
\textbf{Chapter~\ref{chap:conceptual_study}} presents the conceptual study of the platform, modelling
the system through a general use case diagram and actor-specific refinements
covering all four user roles, a class diagram formalizing the domain model,
and sequence diagrams capturing the dynamic behaviour of critical workflows
including VM session provisioning, camera control, and user
authentication. \textbf{Chapter~\ref{chap:realization}} covers the realization and experimentation
phase, describing the software architecture adopted for the platform - its
layered organization, asynchronous provisioning pipeline, and external
infrastructure boundary - presenting the hardware and software development
environment, the principal interfaces implemented for each user role, the
network topology and segmentation strategy, the production deployment and
security configuration, and the verification activities confirming functional
correctness, role-based access control, remote laboratory workflows, and
performance. The report closes with a critical assessment of achievements and
limitations encountered during development, and concrete recommendations for
future platform enhancements.

\cleardoublepage